\documentclass[aps,prb,reprint,amsmath,amssymb,longbibliography,floatfix]{revtex4-2}
\usepackage{booktabs}
\usepackage[colorlinks=true,linkcolor=blue,citecolor=blue,urlcolor=blue]{hyperref}
\hypersetup{pdftitle={Comment on Ergotropy and entanglement in critical spin chains},
 pdfauthor={Xiaojian Zeng},
 pdfsubject={Finite-chain correction to occupation-aligned passive energy},
 pdfkeywords={ergotropy, passivity, fermion chain}}
\newcommand{\Tr}{\operatorname{Tr}}
\newcommand{\ket}[1]{\lvert #1\rangle}
\newcommand{\bra}[1]{\langle #1\rvert}
\newcommand{\al}{\mathrm{al}}
\newcommand{\pa}{\mathrm{p}}
\begin{document}
\title{Comment on ``Ergotropy and entanglement in critical spin chains''}
\author{Xiaojian Zeng}
\email{xiaojianzeng@foxmail.com}
\affiliation{PBC School of Finance, Tsinghua University, Beijing 100190, China}
\date{October 13, 2026}
\begin{abstract}
The occupation-alignment expression in Eq.~(34) of B.~Mula et al.,
Phys. Rev. B \textbf{107}, 075116 (2023), minimizes the subsystem energy under
number-preserving Gaussian orbital rotations. It need not give the minimum
over all subsystem unitaries. For the unchanged half-filled ground state of
the homogeneous open fermion chain, we establish the first failure among
lengths $N=4m$ at $N=12$. A swap between two physical two-particle states
lowers the aligned energy while preserving particle number and parity.
Rational interval calculations enclose the full additional extractable work
between $1.356396207356\times10^{-6}J$ and
$1.356396207357\times10^{-6}J$. We give a quartic extraction gate and finite
comparisons through $N=24$. The corrections amount to about 15--31 parts per
million of the aligned bound energy. These finite results leave the proposed
thermodynamic entropy relation undecided.
\end{abstract}
\maketitle

\section{The optimization being corrected}
Subsystem ergotropy distinguishes the energy removable by a unitary from the
energy remaining above the isolated subsystem ground state. Reference~\cite{Mula2023}
defines the passive energy by ordering the complete many-body spectra in
Eqs.~(1)--(4), but identifies it with an occupation-aligned one-body sum in
Eq.~(34). The distinction matters even for its specified free-fermion
ground-state marginal. Ordering individual occupations against one-particle
energies need not order their many-body products against subset energies.

The general distinction between unrestricted and Gaussian work extraction is
established prior work~\cite{OnumaKalu2018}. Population swaps are also standard.
The result here concerns the actual marginal of the chain in
Ref.~\cite{Mula2023}, without adjustable mode temperatures or a replacement
state. The nearby long-range Kitaev study~\cite{Mitra2025} treats pairing,
chemical potential, and asymptotic entanglement spectra; we make no correction
claim about that model. Our counterexample concerns the exact finite
optimization in Eq.~(34) of Ref.~\cite{Mula2023}. It does not settle the separate
large-$N$ bound-energy prediction in that reference.

\section{Exact input and two operation sets}
Let $N=4m$, $m\geq1$, and $\ell=N/2$. With canonical fermions and $J>0$,
take the open Hamiltonian
\begin{equation}
 H_N=-J\sum_{j=1}^{N-1}(c_j^\dagger c_{j+1}+c_{j+1}^\dagger c_j).
 \label{eq:H}
\end{equation}
There is no chemical potential, pairing, interaction, or bath. The empty
vacuum has zero energy. Let $\rho_A$ be the marginal of the unique half-filled
ground state on $A=\{1,\ldots,\ell\}$, and let $H_A$ retain only the internal
links of $A$ after isolation. Work below is the decrease of $\Tr(\rho_AH_A)$;
it does not include the cost of cutting the chain or implementing a control.

Solving the open one-body difference equation, with vanishing wave function
at sites $0$ and $N+1$, gives sine modes and energies
$-2J\cos[k\pi/(N+1)]$. Filling its $N/2$ negative energies gives exactly
\begin{equation}
 (C_A)_{ij}=\frac{2}{N+1}\sum_{k=1}^{N/2}
 \sin\frac{ik\pi}{N+1}\sin\frac{jk\pi}{N+1},
 \label{eq:C}
\end{equation}
where $(C_A)_{ij}=\langle c_i^\dagger c_j\rangle$, $1\leq i,j\leq\ell$.
Let $\nu_1\geq\cdots\geq\nu_\ell$ be its eigenvalues. The physical energies
of $H_A$, in ascending order, are
\begin{equation}
 \varepsilon_j=-2J\cos\frac{j\pi}{\ell+1},\qquad
 E_{\al}=\sum_{j=1}^{\ell}\nu_j\varepsilon_j.
 \label{eq:align}
\end{equation}
This pairs descending occupations with ascending physical energies, giving the
source expression its favorable convention. Rephasing sites or reversing
one-body labels is not the counterexample.

The marginal is a number-conserving Gaussian state. In its natural orbitals,
Wick factorization fixes all occupation moments and therefore the density
matrix as a product of Bernoulli mode densities. In particular, all its
$2^\ell$ eigenvalues are
\begin{equation}
 q(\mathbf n)=\prod_{j=1}^{\ell}\nu_j^{n_j}(1-\nu_j)^{1-n_j},
 \quad\mathbf n\in\{0,1\}^{\ell}.
 \label{eq:q}
\end{equation}
Let $V$ contain real natural orbitals in descending occupation order, and
$\Phi_{ij}=\sqrt{2/(\ell+1)}\sin[ij\pi/(\ell+1)]$ contain the physical
energy orbitals. The number-preserving Fock lift $\Gamma(W)$ of
$W=\Phi V^\dagger$ maps $\rho_A$ to a state $\rho_{\al}$ whose probabilities
are Eq.~\eqref{eq:q} in the physical-mode occupation basis. Its energy is
$E_{\al}$. The one-body trace rearrangement inequality proves that this is
the minimum over number-preserving Gaussian orbital rotations.

For arbitrary subsystem unitaries, instead let $p_r^\downarrow$ be all
$q(\mathbf n)$ in descending order and $E_r^\uparrow$ all
$E(\mathbf n)=\sum_j n_j\varepsilon_j$ in ascending order. Then
\begin{equation}
 E_{\pa}=\min_U\Tr(U\rho_AU^\dagger H_A)
       =\sum_{r=1}^{2^\ell}p_r^\downarrow E_r^\uparrow.
 \label{eq:passive}
\end{equation}
Indeed, the overlaps between eigenbases form a doubly stochastic matrix;
the trace is a linear function of this matrix, minimized by the opposite
ordering permutation, which is itself unitary. Degeneracies are retained.
Define $\Delta W=E_{\al}-E_{\pa}\geq0$. If operations must commute with
$\widehat n_A=\sum_jc_j^\dagger c_j$, one sorts within each number sector
instead, obtaining $E_{\pa}^{(n)}$ and
$\Delta W_n=E_{\al}-E_{\pa}^{(n)}$. Sector probabilities are neither
discarded nor renormalized.

\section{A charge-conserving counterexample at \texorpdfstring{$N=12$}{N=12}}
For $\ell=6$, let $d_j^\dagger=\sum_i\Phi_{ij}c_i^\dagger$ and choose
\begin{equation}
 \ket a=d_2^\dagger d_3^\dagger\ket{0},\qquad
 \ket b=d_1^\dagger d_4^\dagger\ket{0}.
 \label{eq:ab}
\end{equation}
Both states contain two particles and have even parity, with the other modes
empty. Their energies and aligned probabilities obey the strict rational
enclosures
\begin{align}
 0.335125603737&<\frac{E_b-E_a}{J}<0.335125603738,\nonumber\\
 2.002576633338&<10^6(q_b-q_a)<2.002576633339.
 \label{eq:gaps}
\end{align}
Thus the higher-energy state has higher probability. These signs use the
actual projector~\eqref{eq:C}, not an asymptotic spectrum. For orientation,
the descending occupations are approximately
\begin{equation}
 \begin{split}
 (\nu_1,\ldots,\nu_6)\simeq(&0.999999987458,\ 0.999800520240,\\
 &0.899541716786,\ 0.100458283214,\\
 &0.000199479760,\ 0.000000012542).
 \end{split}
\end{equation}
The approximate entries are not used to decide a sign.

Consider the exact rank-state swap
\begin{equation}
 S=I-\ket a\bra a-\ket b\bra b+\ket a\bra b+\ket b\bra a.
 \label{eq:S}
\end{equation}
It satisfies $S^\dagger=S$, $S^2=I$, and commutes with both $\widehat n_A$
and $(-1)^{\widehat n_A}$. Its work gain from $\rho_{\al}$ is
\begin{equation}
 g=\Tr[(\rho_{\al}-S\rho_{\al}S^\dagger)H_A]
   =(E_b-E_a)(q_b-q_a),
\end{equation}
with
\begin{equation}
 6.71114703278\times10^{-7}<g/J<6.71114703279\times10^{-7}.
 \label{eq:gain}
\end{equation}
Consequently $U=S\Gamma(W)$ takes the original marginal to an energy below
$E_{\al}$. Every density eigenvalue is preserved. The universal finite
equality $E_{\pa}=E_{\al}$ is therefore false, including when particle
number and parity conservation are imposed. No environment operation or
measurement enters this witness.

A quartic operation avoids projectors onto empty spectator modes. Set
\begin{equation}
 K=d_1^\dagger d_4^\dagger d_3d_2+\mathrm{h.c.},\qquad
 U_K=\exp(-i\pi K/2).
\end{equation}
After alignment it exchanges the same two four-mode patterns for every
occupation of modes 5 and 6. Fermionic phases do not affect the energy.
Spectator probabilities sum to one, giving
\begin{equation}
 \begin{split}
 g_K={}&(\varepsilon_1+\varepsilon_4-\varepsilon_2-\varepsilon_3)\\
 &\times[\nu_1(1-\nu_2)(1-\nu_3)\nu_4\\
 &\hspace{1.2em}-(1-\nu_1)\nu_2\nu_3(1-\nu_4)],
 \end{split}
\end{equation}
and
\begin{equation}
 6.71248612208\times10^{-7}<g_K/J<6.71248612209\times10^{-7}.
\end{equation}
This generator conserves number and parity. Its physical-mode interaction
is not asserted to be spatially local, fast, or inexpensive.

\section{First failing size and finite magnitude}
The earlier allowed sizes $N=4,8$ are exactly passive after alignment.
To see this, the staggered matrix $D_{jj}=(-1)^j$ anticommutes with the
one-body hopping. With no zero mode, $DP_-D=I-P_-$ for the occupied
projector. Restriction gives $D_AC_AD_A=I-C_A$, so occupations come in
complementary pairs. Relative to the filled negative-energy local state,
each pair has a hole and a particle excitation with equal positive energy
cost and equal entanglement cost $\kappa=\log[\nu/(1-\nu)]$, $\nu>1/2$.
Probability is proportional to the exponential of minus the total
entanglement cost.

For $N=4$ there is one cost repeated twice. Increasing excitation count
increases energy and decreases probability, proving passivity. For $N=8$
the two energy costs, each repeated twice, are
$J\alpha=J(1+\sqrt5)/2$ and $J\beta=J(\sqrt5-1)/2$, with
$\alpha>2\beta$. Rational interval evaluation of Eq.~\eqref{eq:C} proves
\begin{equation}
 \nu_1(1-\nu_2)^2>(1-\nu_1)\nu_2^2,
\end{equation}
hence $\kappa_1>2\kappa_2$. Excitation counts
$(s_1,s_2)\in\{0,1,2\}^2$ are therefore lexicographically ordered by
both physical and entanglement costs: increasing $s_1$ dominates every
possible decrease in $s_2$, and at fixed $s_1$ both costs increase with
$s_2$. States with the same counts have equal energy and probability.
This proves $E_{\pa}=E_{\al}$ at $N=8$, rather than inferring equality
from a small numerical residual. Thus $N=12$ is the first failure in
the sequence $N=4m$.

\begin{table*}[t]
\caption{Complete finite-spectrum corrections as strict enclosures. The
last column gives $10^6\Delta W/Q_{\al}$ approximately, with
$Q_{\al}=E_{\al}-E_{A,0}$. Both corrections vanish exactly at $N=4,8$.}
\label{tab:work}
\begin{ruledtabular}
\begin{tabular}{c c c r}
$N$ & $10^6\Delta W/J$ & $10^6\Delta W_n/J$ & ppm\\
\hline
12 & $(1.356396207356,\ 1.356396207357)$ & $(1.351749675589,\ 1.351749675590)$ & 15.09\\
16 & $(1.871176478412,\ 1.871176478413)$ & $(1.863882069329,\ 1.863882069330)$ & 23.52\\
20 & $(2.029342670961,\ 2.029342670962)$ & $(2.028803336503,\ 2.028803336504)$ & 28.47\\
24 & $(2.003459793533,\ 2.003459793534)$ & $(2.001606649667,\ 2.001606649668)$ & 31.00\\
\end{tabular}
\end{ruledtabular}
\end{table*}

Table~\ref{tab:work} uses every Fock state at each size, totaling 5460
states over $N=4,8,12,16,20,24$. At $N=12$,
$E_{\al}\simeq-3.40404538383307J$ and
$E_{\pa}\simeq-3.40404674022927J$.
The aligned bound energy is approximately $0.0899138J$, and its aligned
ergotropy is $0.0365547J$. The correction is about 15.09 ppm of the
former and 37.11 ppm of the latter. It is an exact optimization distinction
with a small quantitative effect.

\section{Error control and interpretation}
The strict bounds are computer-assisted enclosures of the exact finite
input. Every arithmetic operation rounds outward on rational intervals
with denominator $10^{90}$. Machin's identity
$\pi=16\arctan(1/5)-4\arctan(1/239)$ is enclosed using alternating
series with 120 and 40 terms and next-term bounds. Angle reduction puts
sines in $[0,\pi/2]$. A degree-199 Taylor polynomial, with its zero
degree-200 coefficient, has remainder bounded by $2^{201}/201!$.
This encloses Eq.~\eqref{eq:C} and all physical energies.

Numerical proposals for eigenvectors are only checked inputs. For every
rational vector $v$ and center $\mu$, exact interval arithmetic verifies
\begin{equation}
 \|(C_A-\mu I)v\|^2<10^{-90}\|v\|^2.
\end{equation}
The Hermitian spectral decomposition then puts an eigenvalue within
$10^{-45}$ of $\mu$. All $\ell$ intervals are strictly disjoint and inside
$(0,1)$, so together they isolate the entire correlation spectrum.
No assumed eigenvector orthogonality or floating eigensolver accuracy
enters this argument.

Every product probability and subset energy is enclosed. Order statistics
are monotone under componentwise bounds; sorting all lower endpoints and
all upper endpoints therefore encloses the corresponding sorted exact
values. Interval products and sums enclose Eq.~\eqref{eq:passive}, including
negative energies and degeneracies, and the same procedure applies within
number or parity blocks. A separate direct ground-state calculation in
the 924-dimensional half-filled $N=12$ sector, followed by an ordinary
site-Fock partial trace, corroborates the marginal, orbital alignment,
swap, and quartic energy traces to better than $2.2\times10^{-15}J$ for the
two gain formulas. This floating cross-check is not the proof of the
strict bounds. The exact checker and its input proposals are supplied
as a reproducibility resource.

The physical correction is
$W=W_{\al}+\Delta W$ and $Q=Q_{\al}-\Delta W$, where
$W_{\al}=\Tr(\rho_AH_A)-E_{\al}$ and $Q=E_{\pa}-E_{A,0}$.
Occupation alignment remains valid with an explicit restriction to
number-preserving Gaussian orbital rotations, or as an upper bound on
the unrestricted passive energy. The additional work survives
number and parity conservation. We establish no classification for
$N>24$, no thermodynamic error estimate, and no experimental or
finite-time protocol. In particular, a finite failure of Eq.~(34) does
not determine the fate of the asymptotic entropy relation in
Ref.~\cite{Mula2023}.

\paragraph{Research tools.}
OpenAI Codex was used for the witness search, derivations,
verification-code development and execution, literature comparison, and
manuscript drafting. Exact arithmetic checks and separate numerical
representations are described above; they do not constitute completed
human scientific review.

\paragraph{Data availability.}
The accompanying reproducibility resource contains the exact checker,
rational eigenvector proposals, complete certificate, and optional
direct-density code and data.

\begin{acknowledgments}
OpenAI Codex assisted with drafting. Human scientific review, authorship
confirmation, and final text approval remain pending.
\end{acknowledgments}

\bibliography{references}
\end{document}
